% Submission target: 2nd Workshop on Grounded and Faithful Vision-Language Models
% for Real-World Deployment (VLM4RWD), NeurIPS 2026. Eight pages excluding
% references and appendices, double-blind, submitted through OpenReview.
%
% Drop neurips_2026.sty beside this file and switch the two lines below to
%   \documentclass{article}\usepackage[final]{neurips_2026}
% once the workshop publishes the style file. Nothing else changes: the geometry
% package here only stands in for that layout.
\documentclass[10pt]{article}
\usepackage[margin=1in]{geometry}
\usepackage{booktabs}
\usepackage{graphicx}
\usepackage{hyperref}
\usepackage{xcolor}
\newcommand{\TBD}[1]{\textcolor{red}{[TBD: #1]}}

% Figures and the main table are build products of the frozen run, never hand
% authored. Copy the outputs of `canyonbench trace report --dataset-dir <dataset>`
% into this directory (flat, beside main.tex: some engines will not resolve
% graphics from subdirectories). Until an artifact exists the float degrades to a
% visible TBD instead of breaking the build, so the draft always compiles from a
% clean checkout and no float can ever be silently empty.
\newcommand{\Generated}[2]{%
  \IfFileExists{#1}{\includegraphics[width=\linewidth]{#1}}%
  {\TBD{generate \texttt{\detokenize{#1}} (#2)}}%
}
\newcommand{\GeneratedTable}[2]{%
  \IfFileExists{#1}{\resizebox{\linewidth}{!}{\input{#1}}}%
  {\TBD{generate \texttt{\detokenize{#1}} (#2)}}%
}

\title{CanyonBench-Trace: Measuring Whether Aerial Vision--Language Models
Depend on the Evidence They Cite}
\author{Anonymous Authors}
\date{}

\begin{document}
\maketitle
\begin{abstract}
A vision--language model deployed on an aerial platform can answer correctly, point
at a plausible region, and still not be using the pixels it cites. Accuracy cannot
separate these cases, and neither can localization overlap. We introduce
CanyonBench-Trace, a procedural benchmark that measures whether an answer
\emph{causally depends} on the visual evidence the model itself names. A
deterministic virtual camera over orthoimagery emits each view together with
pixel-exact feature masks and camera geometry, so counterfactual interventions can
be applied to the true feature, to a covariate-matched irrelevant region, and to
the cells the model verbalizes---using only repeated black-box queries, so
proprietary and open-weight systems are evaluated identically. Because the camera
fixes ground sample distance, the benchmark also contains an optically grounded
hallucination test: features that are mapped and confirmed at low altitude become
physically unresolvable at float altitude in the same scene, so a positive answer
there is ungrounded by imaging geometry rather than by annotator judgment. We
report presence accuracy, localization, and causal reliance as three distinct
quantities across 120 independent sites and 960 views spanning 3--24\,km, together
with per-model acuity thresholds, selective-risk frontiers, and CAVE, a
training-free black-box wrapper that converts unfaithful positives into
abstentions at a measured coverage and inference-call cost.
\TBD{Insert only frozen headline results and the CAVE reliability/coverage effect.}
\end{abstract}

\section{Introduction}
High-altitude aerial imagery creates a sharp deployment problem for
vision--language models: a road, water body, or cultivated field can be
geographically plausible yet absent or below the image's resolvable scale.
Ordinary accuracy cannot show whether a correct answer used visible evidence,
and an overlapping localization can remain causally decorative.

CanyonBench-Trace places an exact virtual camera over frozen orthoimagery and
applies the same projection to RGB, feature masks, and terrain. It then
intervenes on the true feature, a covariate-matched irrelevant region, the
model's own claimed cells, and random or texture-matched cells. This design
measures both oracle-mask reliance and self-reported evidence faithfulness in a
black-box setting.

\paragraph{Contributions.}
We contribute (1) a reproducible three-group, three-class procedural benchmark
with explicit extinction cases; (2) a multi-operator causal evidence protocol
with matched nuisance controls; (3) result-independent validation instruments
for presentation, edit detectability, operator sensitivity, and suppression
efficacy; and (4) CAVE, a training-free accept-or-abstain wrapper based on
necessity, sufficiency, and nuisance tests. \TBD{Adjust wording only if a frozen
instrument requires claim restriction.}

\section{Related Work}
Broad geospatial VLM evaluations such as GEOBench-VLM and OmniEarth cover
diverse perception, reasoning, grounding, and robustness tasks
\cite{danish2025geobench,fu2026omniearth}. CanyonBench-Trace does not compete on
breadth; it contributes a controlled causal instrument for a narrow deployment
slice.

The closest counterfactual work verifies expert-annotated chest-X-ray regions
through foreground and background editing and evaluates attribution against those
causal regions \cite{xiong2026visual}; that combination of correctness filtering,
foreground deletion, and background-invariance checking is the nearest neighbour
to this work, and we position against it directly rather than around it. Its
companion attribution method requires token log-probabilities, which restricts it
to open-weight systems. Counterfactual edits also expose segmentation
hallucination \cite{li2025hallusegbench}, and whole-image controls and benchmarks
of prior versus visual evidence expose answers that persist without supporting
pixels \cite{zafar2026medical,sun2026prive}. Our distinct object is the model's
own spatial self-report: the protocol intervenes on cells the queried black-box
model names, in addition to exact oracle masks, and needs only repeated queries
and a structured answer, so proprietary and open-weight models are evaluated
identically. That prior work's own ablation, in which bounding-box zero masking
outperformed diffusion inpainting for causal attribution, independently supports
our choice of transparent suppression as the primary operator.

ERASER established comprehensiveness, sufficiency, and perturbation-curve
evaluation for rationales \cite{deyoung2020eraser}; Jacovi and Goldberg separate
faithfulness from plausibility \cite{jacovi2020faithfulness}. Demonstrations of
unfaithful verbal reasoning further motivate intervention rather than accepting
self-explanation at face value \cite{turpin2023unfaithful,lanham2023faithfulness}.
Finally, ICE shows that conclusions can vary sharply across intervention
operators \cite{basu2026ice}, motivating our required O1--O3 agreement and
grid/$K$ sensitivity analyses.

\section{Dataset and Registration}
Frames are synchronized to operational-flight telemetry through a visually verified anchor and named by elapsed second. Ground and initialization imagery is excluded; the ascent and float phases form the analysis populations. One-Hz frames are distance/perceptually sampled and grouped into trajectory segments and geographic blocks.

Human annotators label six visible features, image conditions, and binary masks of resolvable living green vegetation. Public land-cover, hydrography, and road layers are annotation aids and weak-label comparisons, not per-frame truth.

For registration, correspondences map frame pixels to metric reference coordinates. Homographies are fitted without two or more held-out control points. A frame enters grounding only when held-out RMSE is no greater than one quarter of a 4-by-4 grid-cell ground width. \TBD{Report counts, agreement, reference products, threshold distribution, and exclusions.}


\section{Evaluation Protocol}
Structured probes ask for binary feature presence, vegetation percentage, and vegetation grid cells. False-premise probes compare a leading prompt with an evidence-first prompt. Free descriptions are secondary and use a blinded judge validated against two human annotators with a rule-based fallback.

We report false-positive rate and F1 for presence, signed error and MAE for vegetation, region precision/recall for grounding, and compliance plus mitigation delta for false premises. Confidence intervals resample trajectory segments. Ascent regressions include image-quality controls, ground-truth prevalence, and phase, with covariance clustered by segment. These are associations, not causal altitude effects.


\section{Results}

\begin{table*}[t]
  \centering
  \GeneratedTable{model_summary.tex}{main table from canyonbench trace report}
  \caption{Main results. Balanced accuracy, false-positive rate, area-penalized
  localization F1, OCRS, SEN, SES, EFS, OSG, prompt-prior gap, extinction
  positive rate, acuity threshold, selective risk at fixed coverage, and CAVE
  coverage. Every column is a build product of the frozen run.}
  \label{tab:main}
\end{table*}

\TBD{Insert the frozen Tier A model/class table, format failures, extinction
curves, and site-level intervals.}

\subsection{Oracle and Self-Reported Evidence}

\begin{figure}[t]
  \centering
  \Generated{accuracy_vs_efs.pdf}{task performance against explanation
  faithfulness}
  \caption{Presence accuracy against explanation faithfulness. The informative
  region is any model that is accurate while weakly grounded, which ordinary
  accuracy cannot distinguish.}
  \label{fig:accuracy-vs-efs}
\end{figure}

\TBD{Report OCRS, SEN, SES, EFS, OSG, target/distractor balance, O1--O3 operator
agreement, and V1/V2/V6 before interpreting null or positive effects.}

\subsection{Acuity and Extinction}

\begin{figure}[t]
  \centering
  \Generated{acuity_curves.pdf}{per-model psychometric acuity curves}
  \Generated{extinction_by_altitude.pdf}{extinction-band positive rate across the
  resolution lattice}
  \caption{Per-model acuity and behaviour inside the extinction band. The acuity
  threshold is the apparent width at which detection probability reaches one
  half, and is obtainable for every model regardless of how the causal results
  resolve.}
  \label{fig:acuity}
\end{figure}

\TBD{Report per-model and per-class acuity thresholds, EPR, and the human
extinction-band confirmation rate.}

\subsection{Prompt Priors and CAVE}

\begin{figure}[t]
  \centering
  \Generated{cave_frontier.pdf}{CAVE reliability--coverage--cost frontier}
  \caption{CAVE frontier. Reliability against retained coverage, coloured by the
  inference-call multiplier, reported as a frontier rather than a single
  improvement number.}
  \label{fig:cave}
\end{figure}

\TBD{Report neutral/false-premise/no-image comparisons and the full
reliability--coverage--cost frontier, including necessity/sufficiency/nuisance
ablations.}

\subsection{Sensitivity}
\TBD{Report geographic groups, clean/degraded, grids 4/6/8, K=3/6/10, date/source
gates, detector, road-width, and O4-secondary analyses.}

\section{Discussion and Limitations}
Procedural projection removes manual registration error but does not make map
sources infallible. Source vintage, consensus-layer dependence, orthoimage
artifacts, limited feature classes, synthetic camera geometry, and intervention
distribution shift remain limitations. Exact masks establish where registered
features project; they do not imply that a human or model can resolve them, so
extinction is analyzed separately.

Causal conclusions are operator-dependent unless O1--O3 agree and V2/V6 pass.
CAVE trades reliability for abstention and repeated-call cost; it is not a
guarantee of safe autonomous action. The first release de-scopes the optional
real-flight temporal tier.

\section{Conclusion}
\TBD{Summarize only outcomes supported by the frozen result and instrument
manifests.}


\bibliographystyle{plain}
\bibliography{references}
\appendix
\section{Reproducibility and Secondary Analyses}
\TBD{Insert source/gate tables, camera equations, degradation calibration,
intervention balance, artifact classifier, complete V1--V6 results, prompt text,
model/runtime/cost manifests, mixed-model tables, BH family, audit form, and all
registered ablations.}

\subsection{Compute and reproducibility}
Dataset construction is CPU-bound and separable from inference. Source
acquisition, virtual-camera generation, gate evaluation, intervention rendering,
and all statistical analysis run on CPU nodes; only the self-served models
require VRAM. Open-weight and remote-sensing models are served under vLLM in
bfloat16 on a single 80\,GB accelerator, one model at a time within one session,
and proprietary models are reached over a commercial API. Quantized serving is
inadmissible: it changes model behavior, and a result obtained under reduced
precision would not describe the named model.

Every query is keyed by a content hash of the model, source image, exact prompt,
intervention, and response schema, so runs are append-only, resumable at the
request that failed, and never pay twice for an identical view. The projected
inference requirement is approximately 15 GPU-hours and the full API volume is
approximately 16{,}600 calls; a measured price pilot re-prices the plan from
observed provider token accounting before any production spend. \TBD{Replace
with the frozen wall-clock, GPU-hour, and dollar totals from the run manifests.}

\section{Outcome Matrix}
\TBD{Map pass/fail combinations for acuity, oracle reliance, self-faithfulness,
and CAVE to the corresponding claim restriction. Preserve null and instrument
failure outcomes rather than omitting them.}

\end{document}
